\chapter*{Glosario}
\addcontentsline{toc}{chapter}{Glosario}

% Mantener en orden alfabético. Agregar cada término a medida que aparezca en los capítulos.
\begin{description}
    \item[Aprendizaje federado] Técnica de aprendizaje automático en la que varios dispositivos entrenan un modelo compartido a partir de sus datos locales y envían solo las actualizaciones del modelo, sin transferir los datos crudos.
    \item[Cuantización] Técnica que reduce la precisión numérica de los parámetros de un modelo (por ejemplo, de punto flotante de 32 bits a enteros de 8 bits) para disminuir su tamaño y acelerar la inferencia.
    \item[Estimación de pose] Técnica de visión por computadora que determina la posición de un conjunto de puntos clave del cuerpo o de las manos a partir de una imagen.
    \item[Glosa] Transcripción de una seña mediante una palabra de una lengua oral, utilizada como representación intermedia en la traducción de lenguas de señas.
    \item[Handedness] Lateralidad de una mano detectada, es decir, si se trata de la mano izquierda o de la derecha.
    \item[Inferencia] Uso de un modelo ya entrenado para obtener una predicción a partir de una entrada nueva.
    \item[Landmark] Punto de referencia que representa la posición de una articulación o extremo de la mano dentro de la imagen.
    \item[Pipeline] Secuencia de etapas de procesamiento en la que la salida de cada etapa es la entrada de la siguiente.
    \item[Seña aislada] Seña individual que se reconoce de forma independiente, sin considerar las señas anteriores o posteriores dentro de una oración.
\end{description}
